\chapter{The Science}
\label{chap:the-science}

% EPISODE COVERAGE: Episode 8 -- Science, LearningProcess, SCIENTIFIC,
% Knowledge, ScientificProcess, TRUTH, Gospel, ReligiousProcess, WITNESSED
%
% AUDIENCE NOTE: This is the philosophical pivot. The tower has been built
% through computation and arithmetic. Now it asks: what does it mean to know
% something? The compiler is asked to hold both science and religion as
% formal structures and to distinguish between them without prejudice.
%
% KEY CONCEPT: Science and religion are both admissibility structures.
% Science updates on evidence. Religion does not. The compiler does not
% care which you prefer -- it only asks whether your continuation is
% consistent with the ledger.
%
% TONE: The most philosophically rich chapter. Accessible but uncompromising.
% The satirical names (GUNGAN, Jar Jar, Meesa) have been making a point
% that is now made explicit.

\section{Two Typeclasses}
\label{se:two-typeclasses}

% Science:
%   inductive Science where
%     | repeatable: Science
%     | hypothesis: Science
%     | theory: Science
%
% Gospel:
%   inductive Gospel where
%     | epiphany: Science
%     | state: Science
%
% These are parallel inductive types. Science has three constructors:
% a repeatable observation, a hypothesis (a proposed explanation),
% and a theory (a hypothesis that has survived repeated testing).
%
% Gospel has two constructors: an epiphany (a revealed truth that
% did not arise from observation) and a state (a fixed position).
%
% The structural difference: Science's constructors are ordered by
% evidence accumulation. A theory requires a hypothesis; a hypothesis
% requires a repeatable. Gospel's constructors require nothing prior.
% The epiphany arrives without preparation.
%
% Sidebar: What makes something a theory? In common usage, "theory"
% means a guess. In science, it means the opposite: a hypothesis that
% has survived every attempt to falsify it. A theory is the most
% reliable kind of knowledge. A hypothesis is still on trial.

\section{The Scientific Method as a Typeclass}
\label{se:scientific-method-typeclass}

% LearningProcess:
%   structure LearningProcess ... where
%     galileo: Cult
%     invariant: Science
%     touch_stove?: ...
%
% The touch_stove? function is the predicate for learning from
% experience. A LearningProcess has a Galileo field (a community of
% practice -- those who have agreed to update on evidence) and an
% invariant Science (the current state of knowledge).
%
% SCIENTIFIC wraps this with predictable?: can the current Science
% make predictions about future observations?
%
% Knowledge:
%   inductive Knowledge where
%     | jarjar: Knowledge
%     | ledger: Knowledge
%
% Two kinds of knowledge: jarjar (intuitive, pre-theoretical,
% unreliable but present) and ledger (recorded, ordered, verifiable).
% The Knowledge type is the honest admission that not all knowledge
% is in the ledger. Some of it is jarjar.
%
% ScientificProcess wraps both with learn?: can this process update
% the Knowledge given a new observation?
%
% Sidebar: Karl Popper and falsifiability. Popper argued that a
% scientific claim is one that could, in principle, be false.
% The predictable? predicate is Lean's version of this: a theory
% is scientific only if it generates predictions that could be
% refuted by observation.

\section{What the Compiler Knows About Religion}
\label{se:compiler-religion}

% ReligiousProcess:
%   structure ReligiousProcess ... where
%     the_literature: Gospel
%     pray?: ...
%
% The pray? function is not satirical. It is the predicate for
% adherence: given a Gospel, does this process conform to it?
%
% The compiler treats ReligiousProcess and ScientificProcess
% with identical structural seriousness. Both are processes.
% Both make admissibility judgments. The difference is in how
% they update: SCIENTIFIC updates when predictable? fails;
% ReligiousProcess does not update, because the epiphany is
% not subject to revision.
%
% TRUTH:
%   class TRUTH ... where
%     scientific_process: ScientificProcess
%     martyred?: ... -> ... -> Prop
%
% The martyred? predicate is the equivalence relation on TRUTH:
% two truths are martyred if they were held simultaneously by a
% process that was forced to revise one of them. This is the
% compiler's model of paradigm shift.
%
% Sidebar: Thomas Kuhn and scientific revolutions. Kuhn argued
% that science does not progress by smooth accumulation. It
% advances through crises and revolutions: a new paradigm
% replaces the old, and knowledge that was central becomes
% peripheral or false. The martyred? predicate formalizes this:
% the old truth was held, the revision was forced, both are recorded.

\section{WITNESSED: The Ledger Closes on Knowledge}
\label{se:witnessed}

% WITNESSED:
%   class WITNESSED ... where
%     baptism: ReligiousProcess
%     witness: Gospel
%     risen?: ...
%
% WITNESSED is the first typeclass that simultaneously requires
% a ReligiousProcess and connects it to the Gospel via a witness.
% The risen? predicate asks: has this Gospel survived the revision
% that the scientific process forced?
%
% This is the tower's model of convergence between science and
% religion. A Gospel that survives the scientific process is not
% refuted -- it has been witnessed. A Gospel that does not survive
% is not necessarily false -- it has not been witnessed within
% this ledger.
%
% The compiler makes no judgment about which Gospels are true.
% It only asks: is this process consistent? Has the witness been
% recorded? Is the continuation admissible?
%
% End of chapter: the tower has now agreed to hold Science and
% Gospel as first-class values, to distinguish knowledge that
% updates from knowledge that does not, and to record the
% history of revisions. The ledger is now epistemological.
% The next question is: what does the universe look like from
% inside a ledger like this?

